\documentclass[10pt,a4paper]{amsart}
\usepackage[margin=19mm]{geometry}
\usepackage{amsmath,amssymb,mathtools}
\usepackage[scr=rsfs,cal=euler]{mathalfa}
\usepackage{xfrac}
\usepackage[unicode,hidelinks,bookmarks=false]{hyperref}
\newcommand{\ie}{i.e.\ }
\newcommand{\cf}{cf.\ }
\newcommand{\resp}{respectively\ }
\newcommand{\Subsection}{\subsection}
\providecommand{\nobreakdash}{\nobreak\hbox{-}}
\setlength{\emergencystretch}{2em}
\multlinegap0pt
\usepackage{amssymb}
\usepackage{mathtools}
\usepackage[shortlabels]{enumitem}
\usepackage{xcolor}
\usepackage{tikz}
\newtheorem{lemma}{Lemma}
\title{Torsors under the universal Jacobian over $\mathcal{M}_g$}
\author{Qixiao MA}
\date{Source: arXiv:2607.17206v2 (CC BY 4.0)}
\begin{document}
\maketitle

\noindent\emph{Source and licence.} Torsors under the universal Jacobian over $\mathcal{M}_g$. Version arXiv:2607.17206v2 (CC BY 4.0), distributed by arXiv under the Creative Commons Attribution 4.0 International licence, http://creativecommons.org/licenses/by/4.0/, as recorded in the arXiv licence metadata for this exact version and shown on the abstract page https://arxiv.org/abs/2607.17206v2. Full licence notice: \texttt{licenses/arxiv-2607.17206-CC-BY-4.0.txt}.\par\smallskip

\noindent\textit{Editorial scope.} Counted unit: the article's lemma identify (source0.tex lines 174--175). The article's complete original proof of it is delivered with it (source0.tex lines 176--177, one \texttt{proof} environment of 2 lines). No mathematics is written, repaired or re-derived: the statement and the proof are the article's own lines, in the article's own notation, and no retained line is substituted or re-flowed.  No further source-local context is needed: every cross-reference of every retained line resolves inside the two delivered spans.  Nothing was omitted from any retained span, so the package carries no omission record.  No retained line contains a citation, so the wrapper carries no bibliography and no bibliography entry.  No retained line contains a figure environment, an image-inclusion command or drawing code, so no image is delivered and the package declares no asset dependency.  Editorial formatting only, in addition: an AMS \texttt{amsart} preamble that restates the article's own declarations and macros used by the retained lines, namely the environments \texttt{lemma} on separate counters, together with the standard packages those lines need.  The retained environments print in this document as Lemma 1 in order of appearance, whereas the article prints the same items with the numbering of its own full document.  The complete native source of the article as distributed in the arXiv e-print for this exact version is bundled unchanged as \texttt{sources/205545/source0.tex}.
\begin{lemma}\label{identify}Under the identification $\mathrm{H}^2(\Gamma_g,\mathbb{H})\cong\mathrm{H}^2(\mathcal{M}_g,\mathrm{R}^1\pi_*\mathbb{Z})$, up to a $\pm$ sign which depends on convention, this class satisfies $\epsilon(1)=\mathrm{d}_{2,\mathbb{Z}}^{0,1}(1)$, where $\mathrm{d}_{2,\mathbb{Z}}^{0,1}\colon\mathrm{H}^0(\mathcal{M}_g,\mathrm{R}^2\pi_*\mathbb{Z})\to\mathrm{H}^2(\mathcal{M}_g,\mathrm{R}^1\pi_*\mathbb{Z})$ is the transgression map.
\end{lemma}

\begin{proof}The transgression map can be calculated by \v{C}ech cohomology with respect to the Teichm\"uller covering $\{\mathcal{T}_{g}\to\mathcal{M}_{g}\}$. The \v{C}ech nerve $T_{g}^{\times_{\mathcal{M}_g}k+1}\cong(\Gamma_g)^k\times\mathcal{T}_g$ and the \v{C}ech complex is isomorphic to the bar complex for group cohomology of $\Gamma_g$.
\end{proof}

\end{document}
